% -*- TeX:UK -*-
\documentclass[10pt]{beamer}
\usetheme{metropolis}
%\useinnertheme{rectangles}
\setbeamercovered{%
still covered={\opaqueness<1->{15}},
again covered={\opaqueness<1->{40}}}

\hypersetup{colorlinks,linkcolor=black,urlcolor=brown,citecolor=brown}

\include{shared-tex/set-up-fonts-and-icons}

\begin{document}

\title{PBIO-141\\Sensory and Physiological Ecology\\of  Plants}
\subtitle{2: Terminology and viewpoint}
\author{Pedro J. Aphalo}
\date{January--February 2022}
\institute[Univ.\ of Helsinki]{M.Sc.\ in Plant Biology, University of Helsinki\\[2ex] \url{http://blogs.helsinki.fi/aphalo/}\\[2ex] \url{mailto:pedro.aphalo@helsinki.fi}}

\begin{frame}
\maketitle
\end{frame}

  \begin{frame}[c]
    \begin{center}
      \begin{small}
        \copyright 2006--2022 by Pedro J. Aphalo\\
        University of Helsinki, Finland.\\
        \textcolor{blue}{\url{http://blogs.helsinki.fi/senpep-blog/}}\\[2ex]
      \end{small}

      \begin{footnotesize}
        Sensory and Physiological Ecology of Plants slides by Pedro J. Aphalo are licensed under a Creative Commons Attribution-ShareAlike 4.0 International License.

      \includegraphics[width=6em]{../figures/copyright/by-sa}\\[2ex]
      \end{footnotesize}
        
        \begin{scriptsize}
        Typeset in Lucida Sans, \textrm{Luicda Bright}, \texttt{Lucida Console} and Lucida Math. Icons from fonts ``WebHostingHub Glyphs'' (under SIL-Open Font License) from \url{https://www.webhostinghub.com/}; ``insect icons'' (free from \url{http://www.woodcutter.es/}); ``linea-basic-10'' and ``linea-weather-10'' (free from \url{https://github.com/linea-io}), ``Mini Pics Uprooted Twig'' and ``Mini Pics Uprooted Twig'' (commercial, from Image Club Graphics, Inc.). Plant icon as .svg by Abdul Wahhab (free from \url{NounProject.com}).

        Illustrations and text quoted from copyrighted sources is excluded from this license and their use should respect the original licenses.
        \end{scriptsize}
    \end{center}
  \end{frame}


\begin{frame}
    \frametitle{Outline}
    \tableofcontents
\end{frame}

\section{About the course}

\begin{frame}{Learning objectives (in brief)}
    Are that at the end of this course:
    \begin{itemize}
    \item Learn the language of sensory and physiological ecology of plants.
    \item Learn the basic "workings" of plants and how they have adapted and how they can acclimate to environment.
    \item Learn concepts and a way of looking at plants based on the current scientific ideas.
    \item Learn how to find answers to questions about acclimation and adaptation of plants through experiments.
    \item Learn how to find, interpret and use information to solve problems and answer questions related to plant ecology.
    \item However, you will \textbf{not} have to memorize detailed information that is available in books, journals or the internet.
    \end{itemize}
\end{frame}

\begin{frame}{Contents of the course}
More details available in Moodle. All the points listed below are discussed emphasising acclimation and adaptation.
\begin{itemize}
   \item Introduction
   \item The shared environment of plants
   \item Sensory biology
   \item Adaptation and acclimation of carbon and energy metabolism
   \item Adaptation and acclimation of water uptake and loss
   \item Adaptation and acclimation to mineral nutrient supply
   \item Adaption and acclimation of morphology and anatomy
   \item Plasticity in a changing or artificial habitat
\end{itemize}
\end{frame}

\begin{frame}{Understanding vs.\ memorization}

  \begin{alertblock}{Teaching and learning}
    \begin{itemize}
      \item The role of examples (in this course and elsewhere) is to make a concept tangible\ldots
      \item \ldots easier to grasp than a purely abstract idea.
      \item When studying focus on the concepts\ldots
      \item \ldots once you grasp a concept, or learn the idea\ldots
      \item \ldots you will be able to find or make up another suitable example when needed.
    \end{itemize}
  \end{alertblock}

\end{frame}

\begin{frame}{Examples}

\begin{itemize}
  \item I will use examples from both wild species and cultivated species. Mostly higher plants from terrestrial environments.
  \item In recent years physiological research has made frequent use of \emph{Arabidopsis thaliana}.
  \item \emph{Arabidopsis thaliana} is a useful model, but the variation in the plant kingdom is enormous: in size, shape, function, growth rate, and lifespan.
  \item This variation is the result, at least in part, of variation in past and current growth environments.
  \item Natural selection is constrained by evolutionary history.
  \item Random events also play a role, especially in small populations.
\end{itemize}

\end{frame}\section{General system theory}
\nocite{Meadows2008,Jones2015}

\begin{frame}{Basic concepts}
  \begin{itemize}
    \item Similar structures and properties are present in all systems.
    \item Possible behaviour is determined by the structure of a ``whole thing''.
    \item The structure emerges from the interactions among its ``component parts''.
    \item If interactions exist, then the contribution of the parts to the behaviour of the ``whole'' depends on the nature of interactions.
    \item Behaviour of a complex whole cannot be predicted from the behaviour of the parts.
    \item ``Emergent properties'' appear as a result of the interactions.
  \end{itemize}
\end{frame}

\begin{frame}{Basic concepts}
  \begin{itemize}
    \item Systems can be nested in other systems.
    \item Thus a hierarchy appears.
    \item A crucial concept is feedback, which can be positive or negative.
    \item \emph{Thinking in systems} is applicable to engineering, computing, human society, business administration, biology, ecology, medicine, etc.
    \item Even to managing everyday life by an individual.
    \item At its simplest, \emph{Thinking in systems} .can be defined as being aware of the ``big picture''.
  \end{itemize}
\end{frame}

\section{Hierarchy and scale}
\nocite{Allen1982}
\begin{frame}{What is hierarchy? \Discussion 5 + 5 min}
  \begin{alertblock}{Hierarchies in biology \Discussion}
    \begin{itemize}
      \item Think examples of the application of the concept of \emph{hierarchy} in biology.
    \end{itemize}
  \end{alertblock}
\end{frame}

\begin{frame}{Hierarchy I}

The term \textit{hierarchy} is frequently used in an intuitive way,
but Allen and Starr (1982) have developed a theory, mainly focused
on ecological systems.

\begin{quotation}
By \textit{hierarchy} we understand a system of behaviour
relationships where upper levels limit and control the lower levels
in a greater or smaller measure depending on the \textbf{time
constants} of their behaviour.
\end{quotation}

\end{frame}

\begin{frame}{Hierarchy II}
\begin{itemize}
\item Where \textit{time constant} of a response is the length of time
between when a response to a stimulus starts to be observable until
it reaches a certain level close to the maximum that the given
stimulus will induce. It is a measure of the speed of change of a
response to an instantaneous stimulus.

\item In contrast \textit{lag} is the delay between the application of
the stimulus and the start of an observable response.
\end{itemize}
\end{frame}

\begin{frame}{Hierarchy III}

The degree of control, that is the asymmetry of the relationship,
depends on the time constants of the behaviour.

\vspace{\stretch{1}} \center{\fbox{\parbox{0.8\textwidth}{When we
study the mechanism of a system at one level, we can ignore what changes much more slowly
(has a larger time constant).}}} \vspace{\stretch{1}}

\vspace{\stretch{1}} \DExamples\vspace{\stretch{1}}

\end{frame}

\begin{frame}{Hierarchy IV: examples}
\begin{itemize}
\item For example if we study the evolution of plants on Earth we can
ignore the big bang and the expansion of the universe (and many
other things).

\item If we are interested \emph{only} (narrowly?) in the response of photosynthesis to light
in Arabidopsis, we can ignore the evolution of plant species.

\item If we are interested only in the primary reactions of photosynthesis
we can use in our experiments isolated chloroplasts and ignore the
plant.
\end{itemize}
\end{frame}

\begin{frame}{Levels of organization}
An example of a hierarchy is the hierarchy of levels of organization

Ecosystem $\rightarrow$\\
\hspace{.5cm}Community $\rightarrow$\\
\hspace{1cm}Population $\rightarrow$\\
\hspace{1.5cm}Individual $\rightarrow$\\
\hspace{2cm}Organ $\rightarrow$\\
\hspace{2.5cm}Tissue $\rightarrow$\\
\hspace{3cm}Cell $\rightarrow$\\
\hspace{3.5cm}Organelle $\rightarrow$\\
\hspace{4cm}Molecule $\rightarrow$\\
\hspace{4.5cm}Atom $\rightarrow$\\
\hspace{5cm}Subatomic particle

\end{frame}

\begin{frame}{Scale}
    \begin{description}[type=1]
        \item[Temporal scale] Fast physiological
        responses, acclimation, and adaptation happen at
        different temporal scales. The processes involved are
        different.
        \item[Spatial scale] For example when
        studying transpiration, we can do it at different
        scales: the leaf, the tree crown, the forest, etc.
        The main controlling environmental variables may
        be different.
    \end{description}
\end{frame}

\begin{frame}{Mechanism and `reason' (=justification)}
    \begin{itemize}
        \item \emph{Upward and downward causation} in biology.
        \item We achieve a \emph{mechanistic explanation} for an
        observed phenomenon by studying the levels below.
        e.g. We study whole plant growth in two different habitats.
        We build a mechanistic
        explanation from the growth and responses of individual
        organs: leaves, stems, roots.

        \item We build a \emph{narrative explanation} for the phenomenon
        by looking at the levels above.
        e.g. In this example we find a reason or justification for
        the differences based on natural selection, but we do not study
        the evolutionary process itself.

        \item Of course we can set the focus at any level in the hierarchy and what is mechanism or narrative explanation will move along.~\DExamples
    \end{itemize}
\end{frame}

\section{A video for discussion}
\nocite{Falik2011,Falik2012,Novoplansky2021}

\begin{frame}{Talk by Ariel Novoplansky \Video}

Observation without preconceptions plus theory are important when developing new hypotheses\ldots

\includegraphics[width=0.5\linewidth]{photos/ariel}

 \Video (15 min) \url{https://youtu.be/aClSp71zfro}

\end{frame}


\begin{frame}{Next steps \Discussion}
\begin{itemize}
  \item By now I expect you and me to share enough of a language and frame of mind to discuss without misunderstandings.
  \item If some ideas remain unclear, do ask at any time.
  \item If we had enough time we have watched the video of the TEDx talk by Prof.\ Ariel Novoplansky, but if not please watch it at home. Discuss about it in the corresponding Moodle forum. Remember that discussion can include what implications the phenomena described may have in a question of your own interest, how watching the talk may have changed your views about plants, questions about something that you did not understand, or something that you would like to criticize.
\end{itemize}

\end{frame}

\section*{References}
  \nocite{Ballare2009}
  \begin{frame}[t,allowframebreaks]
    \frametitle{References}
    \printbibliography
  \end{frame}

\end{document}

